% ch07-architecture.tex – Chapter 7: system architecture and innovation (final round, Writer).
\ChapterMission{המשימה: להראות כיצד החישה, הבקר, התצוגה והקמפיין פועלים יחד, ומה חדש בשילוב}
\ChapterIcon{signal}
\chapter{ארכיטקטורת המערכת וחדשנות}\label{chap:architecture}
\begin{scenebox}
לילה. אדם אחד נעצר על המדרכה הריקה. כמה רגעים אחר כך הרמזור לרכב מתחלף לצהוב, והאיש הירוק נדלק.
הוא חוצה בנחת. כשהאור מתחיל להבהב, מופיעות ספרות כתומות ויורדות, והוא מגיע לצד השני לפני שהן נכבות.
\end{scenebox}

הרמזור שבתמונה הבחין בהולך היחיד; הפרק מסביר כיצד, ובאילו גבולות.
הוא מתאר את רכיבי המערכת, את זרימת הנתונים, את לוגיקת האדום המסתגל, את הספירה במופע הפינוי ואת הגנת הפרטיות, ומסיים בחידוש שבשילוב.

\section{מבנה המערכת וזרימת הנתונים}
איור~\ref{fig:t2} מציג את ארבע השכבות. גלאי בשפת המדרכה מזהה ממתינים באזור ההמתנה, סופר אותם ומודד לכל אחד את זמן המתנתו. הבקר מקבל מן הגלאי שלושה נתונים בלבד: מספר הממתינים, זמן ההמתנה הארוך ביותר והאם הממתין יחיד, ומחליט אם ומתי לקצר את האדום. ראשי האות מציגים את המופעים, וראש הספירה מקבל את תזמון הפינוי מן הבקר. שכבת הנתונים שומרת ספירות וזמנים מצטברים לכל מחזור, והם משמשים להערכת הפיילוט ולתצוגת הקמפיין.

\begin{figure}[ht]
  \centering
  \inputfigure{figures/tikz/t2-architecture}
  \caption{ארכיטקטורת המערכת: חישה בשפת המדרכה, בקר עם לוגיקת אדום מסתגל בגבולות בטוחים, ראשי אות עם ספירה במופע הפינוי בלבד, ושכבת קמפיין ונתונים}
  \label{fig:t2}
\end{figure}

\section{לוגיקת האדום המסתגל}
הלוגיקה פועלת בכל מחזור לפי ארבעה כללים:
\begin{enumerate}
\item \textbf{אין ממתינים:} הבקר פועל לפי תוכנית התזמון הקבועה.
\item \textbf{סף המתנה:} להולך יחיד נקבע סף המתנה נמוך מאשר לקבוצה, משום שאין מי שירסן אותו \cite{rosenbloom2009,zhu2021}. ההפרות מצטופפות בתחילת ההמתנה \cite{xing2022}, ולכן הסף קצר; ערכו נקבע בסימולציה.
\item \textbf{בדיקת הגבולות:} בחציית הסף הבקר בודק את הירוק המזערי לרכב, את תקרת הקיצור ואת חלון התיאום (\en{REQ-05}, \en{REQ-08}), ורק אז מסיים מוקדם את הירוק לרכב ופותח מופע הליכה מלא. ההליכה, הפינוי ומרווח הביטחון אינם מקוצרים לעולם (\en{REQ-06}, \en{REQ-07}).
\item \textbf{כשל:} בתקלה בגלאי או בתקשורת הבקר חוזר לתזמון הקבוע (\en{REQ-10}).
\end{enumerate}

\section{ספירה בפינוי ופרטיות}
הספירה מתחילה עם הבהוב האות להולכי הרגל ויורדת עד אפס, ובזמן האדום הראש חשוך \cite{fhwa2025mutcd}. איור~\ref{fig:photo} מראה ראש רמזור מסוג זה בשטח: הספירה מוצגת לצד האות הירוק להולכי הרגל, כלומר בזמן שבו מי שכבר בכביש צריך לדעת כמה זמן נותר לו, ולא בזמן ההמתנה באדום. מכיוון שהאדום המסתגל משנה את מועדי המופעים, הספירה מתוזמנת ישירות מן הבקר (\en{REQ-09}).
הפרטיות מובנית בתכנון: אם הגלאי מצלם, התמונה מעובדת בתוכו ואינה יוצאת ממנו. גלאי מסחרי כזה מסוגל להזרים וידאו \cite{agd2024agd645}, ולכן אפשרות זו מושבתת, ולבקר ולשכבת הנתונים עוברים מספרים בלבד (\en{REQ-11}). גלאי שאינו מצלם ועומד בדרישות הדיוק מועדף.

\authorphoto[72mm]{figures/photo-crosswalk.jpg}{ראש רמזור להולכי רגל המציג ספירה לאחור לצד האות הירוק, מעל מעבר חצייה}{fig:photo}

\section{החידוש}
כל רכיב בנפרד קיים בשוק (פרק~\ref{chap:market}). החידוש הוא בשילוב: גלאים קיימים משמשים לבקשת מופע או להארכת פינוי, ואילו כאן הם משמשים לקיצור האדום לפי מספר הממתינים וזמן המתנתם, בעדיפות להולך היחיד. כך הופך ממצא תצפיתי של מאמר הדגל לכלל בקרה, ושכבת הנתונים מספקת את מדידת ההמתנה שחסרה במאמר \cite{rosenbloom2009}.
